\begin{figure}[!t]
\centering
\includegraphics[width=\linewidth]{exp1_protocol_four_panels.pdf}
\caption{\textbf{Prospective forecasting study.} An agent forecasts
an unresolved question, then receives one generated evidence packet with search
off. The packet's intended direction is hidden from the agent and used only for
scoring. Directional correctness asks whether the update follows a pro- or
anti-$H_1$ packet; selectivity compares movement under directional and
orthogonal packets.}
\label{fig:exp1-pipeline}
\end{figure}

\section{Do forecasts change for the right reason?}
\label{sec:evidence}

We begin on unresolved markets, where the questions are real but lexical cues
remain. We then use Coin City, a controlled three-city forecasting task, to
display competing response regimes and change which regime the context supports
while holding the numerical record fixed. Arbitrary labels remove any stable
meaning from the cue. Across studies, contrasts are paired by
market or episode, and uncertainty resamples that unit. Repeated generations are
not treated as independent observations.

\subsection{Separate evidence from topical chatter}
\label{sec:exp1}

\paragraph{Design.}
We sample 100 unresolved binary Polymarket questions spanning politics,
government, elections, economics, technology, and geopolitics. Each model first
constructs an event model and forecast with web search available. Search is then
disabled, and the model receives one fictional news packet at a time: three
support the YES hypothesis, three support NO, and three concern the same topic
without bearing on the resolution. The model sees only its prior and the snippet.
Packet type, intended direction, and resolution are withheld and used only for
scoring. Nine deployments produce 30,094 valid updates
(Appendix~\ref{app:exp1}).

The design separates two questions. \emph{Selectivity} asks whether the model
changes its forecast more for outcome-bearing evidence than for a topic-matched
control. \emph{Directional consistency} asks whether a sufficiently large change
follows the packet's withheld implication for the market outcome. A
market-held-out bag-of-words classifier measures how much direction is
available from lexical content alone.
Figure~\ref{fig:exp1-pipeline} summarizes the forecast, evidence, and
scoring sequence.


\paragraph{Finding: models first decide whether an update is warranted.}
Figure~\ref{fig:exp1-updating} shows both parts of the result. Every deployment leaves its prior unchanged more often for a topic-matched
control than for directional evidence. When they do update, models move
$1.9$--$7.5\times$ farther for directional packets, with every clustered interval
above one. They also move in the implied direction: directional consistency spans
$.887$--$.990$, compared with $.688$ for the held-out lexical classifier. Eight quality-eligible reviewers match the frozen packet direction on 75.0\%
of judgments, and the panel majority matches 17 of 18 packets.

The first experiment therefore establishes relevance-sensitive updating. Coin
City then tests relationship selection directly by displaying and manipulating
generated response regimes, separating the operative relationship from the
natural-language packets used above.

\begin{figure}[!ht]
\centering
\begin{subfigure}[c]{0.525\textwidth}
\centering
\includegraphics[width=.92\linewidth]{exp1_direction_selective_updating.pdf}
\caption{}
\label{fig:exp1-bars}
\end{subfigure}\hspace{0.01\textwidth}
\begin{subfigure}[c]{0.455\textwidth}
\centering
\footnotesize
\setlength{\tabcolsep}{3pt}
\begin{tabular*}{\linewidth}{@{}l@{\extracolsep{\fill}}rccc@{}}
\toprule
\textbf{Model} & \textbf{Size} & \textbf{Abst.} & \textbf{Sens.}$\uparrow$ & \textbf{EHC$_{\geq3}$}$\uparrow$ \\
\midrule
\multicolumn{5}{@{}l}{\emph{Hosted API (parameters undisclosed)}}\\
\micon{claude.png}~Opus~4.8 & --- & 71.4\% & \cellcolor{heatgreen}$7.0\times$ & \cellcolor{heatgreen}.990 \\
\micon{openai.png}~GPT-5.4 & --- & 8.4\% & \cellcolor{heatgreen}$7.5\times$ & \cellcolor{heatgreen}.987 \\
\micon{deepseek.png}~V4-Pro & --- & 44.4\% & \cellcolor{heatyellow}$4.9\times$ & \cellcolor{heatyellow}.977 \\
\midrule
\multicolumn{5}{@{}l}{\emph{Open-weight, descending parameters} $\downarrow$}\\
\micon{qwen.pdf}~2.5-72B & 72B & 24.2\% & \cellcolor{heatorange}$3.2\times$ & \cellcolor{heatyellow}.963 \\
\micon{llamma.pdf}~3.3-70B & 70B & 10.7\% & \cellcolor{heatyellow}$5.7\times$ & \cellcolor{heatyellow}.965 \\
\micon{llamma.pdf}~3.1-70B & 70B & 2.5\% & \cellcolor{heatyellow}$4.5\times$ & \cellcolor{heatyellow}.961 \\
\micon{qwen.pdf}~2.5-32B & 32B & 11.1\% & \cellcolor{heatorange}$2.8\times$ & \cellcolor{heatyellow}.964 \\
\micon{qwen.pdf}~2.5-14B & 14B & 26.1\% & \cellcolor{heatred}$2.2\times$ & \cellcolor{heatorange}.943 \\
\micon{llamma.pdf}~3.1-8B & 8B & 1.1\% & \cellcolor{heatred}$1.9\times$ & \cellcolor{heatred}.887 \\
\bottomrule
\end{tabular*}
\caption{}
\label{fig:exp1-table}
\label{tab:consistency}
\end{subfigure}
\caption{\textbf{Direction-selective updating across nine models and 30,094
valid updates.} \textbf{(a)} Mean signed revision with 95\%
market-clustered bootstrap intervals. \textbf{(b)} Abstention is the
unchanged-prior rate on topical controls, sensitivity is the conditional
movement ratio, and EHC is directional consistency at a three-point threshold.
Appendix~\ref{app:exp1-threshold-robustness} reports the threshold sweep.}
\label{fig:exp1-updating}
\end{figure}

\subsection{Let context select a relationship, then challenge it}
\label{sec:exp2}

\paragraph{Coin City makes the latent choice observable.}
We call the task Coin City because it carries
Figure~\ref{fig:world-knowledge-evidence}'s unfamiliar-coin problem into city
polling. It turns that ambiguity into a controlled test: two candidate relationships are displayed, context favors one,
and local observations can overturn the initial choice.
Each episode describes three cities whose polls respond to campaign news. Cities
A and B provide noisy examples of a strong and a weak response regime. City C is
the target. A one-sentence description links City C to one reference environment,
and from zero to four City C cases provide increasingly direct local evidence.
For city \(j\) in episode \(i\),
\[
  \Delta p_{ij}=g_j x_{ij}+\epsilon_{ij},
  \qquad \epsilon_{ij}\sim\mathcal{N}(0,5^2),
  \label{eq:coin-city-dgp}
\]
where \(x_{ij}\) is signed net news and \(g_j\) is the city's stable response.
No prompt states the coefficient or names a graph.
Figure~\ref{fig:coin-city-prompt} shows the shared query, reference cities, and
target-city context in one illustrated arm.

\begin{figure}[!b]
\centering
\includegraphics[
  width=.94\linewidth,
  trim=11pt 14.5pt 3.5pt 7pt,
  clip
]{exp2_coin_city_prompt_arms.pdf}
\caption{\textbf{Coin City turns contextual reasoning into a paired
comparison.} Every arm shares the query and numerical cases. The context arm
adds a description connecting City C to one displayed media environment; the
opposite-context arm changes only that sentence. Target cases then accumulate,
allowing local evidence to confirm or correct the contextual choice.}
\label{fig:coin-city-prompt}
\end{figure}

Five prompt arms expose the choice behind the forecast: target cases alone;
references without City C context; references with the correct context; the same
references with the context inverted; and an arbitrary-label version that
replaces familiar media descriptions with KIV and ZOR. In the last arm, which
label names the strong regime changes by episode and must be inferred from the
reference cases. Semantic arms evaluate six hosted deployments over 250 paired
episodes and five evidence depths; the arbitrary-label roster contains fifteen
hosted and open systems. Forecast error is measured against the noise-free
episode expectation, while the implied response coefficient reveals which
reference relationship the forecast follows.

\begin{figure}[!ht]
\centering
\includegraphics[width=\linewidth]{exp2_coin_city_results.pdf}
\vspace{3pt}
\begin{minipage}{\linewidth}
\centering
\footnotesize
\setlength{\tabcolsep}{4pt}
\resizebox{\linewidth}{!}{%
\begin{tabular}{l r rrrr rr}
\toprule
& & \multicolumn{4}{c}{Forecast MAE [95\% CI]} & \multicolumn{2}{c}{$\rho(\hat g,g_C)$} \\
\cmidrule(lr){3-6}\cmidrule(lr){7-8}
Model & Paired $n$ & Target only & No C context & $+$ C context &
$+$ wrong C context & No context & $+$ context \\
\midrule
\micon{claude.png}~Opus~4.8 & 250 & 2.32 [2.20,\,2.46] & 2.65 [2.41,\,2.89] & \textbf{1.82} [1.64,\,2.02] & 4.20 [3.87,\,4.54] & $-.06$ & \textbf{.70} \\
\micon{claude.png}~Opus~5 & 250 & 2.36 [2.12,\,2.59] & 2.62 [2.38,\,2.84] & \textbf{1.76} [1.57,\,1.99] & 4.25 [3.92,\,4.57] & $-.01$ & \textbf{.62} \\
\micon{openai.png}~GPT-5.6 & 250 & 3.60 [3.28,\,3.94] & 2.52 [2.33,\,2.75] & \textbf{1.60} [1.42,\,1.77] & 4.23 [3.90,\,4.55] & $-.04$ & \textbf{.68} \\
\micon{deepseek.png}~V4-Pro & 250 & \textbf{2.45} [2.27,\,2.61] & 2.84 [2.55,\,3.15] & 2.47 [2.19,\,2.77] & 4.36 [4.01,\,4.74] & $-.03$ & \textbf{.52} \\
\micon{kimi.png}~K3 & 250 & 3.35 [3.04,\,3.67] & 2.55 [2.32,\,2.78] & \textbf{1.89} [1.66,\,2.13] & 4.31 [3.98,\,4.62] & $-.02$ & \textbf{.59} \\
\micon{gemini.png}~3.6 Flash & 250 & 2.53 [2.35,\,2.73] & 2.55 [2.31,\,2.80] & \textbf{2.17} [1.89,\,2.48] & 4.29 [3.94,\,4.66] & $.08$ & \textbf{.56} \\
\bottomrule
\end{tabular}
}
\end{minipage}
\caption{\textbf{Context improves zero-shot forecasts and regime assignment
across six deployments.} Panels add references, the correct cue, then an inverted
cue as target observations accumulate. The table reports zero-shot MAE and slope
correlation with 95\% intervals (Appendix~\ref{app:coin-city-results}).}
\label{fig:coin-city-results}
\end{figure}

\begin{wrapfigure}{r}{0.43\textwidth}
\vspace{-1.8em}
\centering
\includegraphics[width=\linewidth]{exp2_symbol_decoding.pdf}
\vspace{-0.8em}
\caption{\textbf{Across model families, larger checkpoints decode the arbitrary
cue, then follow direct evidence.} Lines show group medians and bands span
checkpoints. The split appears within both Qwen3 and Llama~3.1: four larger Qwen
and Llama checkpoints (14--72B) recover the mapping, while three smaller ones
(4--8B) do not. Held-out AUC on 88 held-out episodes.}
\label{fig:coin-city-symbol-decoding}
\vspace{-3.4em}
\end{wrapfigure}

\paragraph{Finding: the forecast follows the selected analogue.}
Figure~\ref{fig:coin-city-results} shows how the contextual cue and target
observations jointly determine the forecast. At zero target cases, correct context lowers error for every deployment; changing
only the context sentence reverses the selected regime. The implied response
loads on the reference named by the cue, not on the cue's words alone:
after partialling out the true regime, the named-reference correlation is
$.51$--$.86$, while the other reference is near zero, and inversion transfers
that loading to the incorrectly named reference. As City C observations
accumulate, the cue's advantage and the harm from inversion shrink. Forecasts
weight context when direct evidence is sparse and shift toward local data as it
becomes more diagnostic.

\subsection{Infer a convention with no stable semantics}
\label{sec:novel-mapping}

The semantic cue still carries a familiar association: national coverage is
described as producing larger movements. KIV and ZOR remove that shortcut. Their
meaning changes across episodes, so a fixed preference for either token is
useless. On the original Coin City population, twelve of fifteen systems recover
the local mapping at zero target cases, including all hosted systems tested. A harder population compresses every score toward chance without
reordering the checkpoints, so the threshold at which recovery is detectable
moves while the ordering holds
(Appendix~\ref{app:coin-city-robustness}).


Figure~\ref{fig:coin-city-symbol-decoding} reports the representation analysis.
Across seven open-weight checkpoints from the Qwen and Llama families, the four
larger checkpoints (14--72B)
recover the arbitrary mapping in their forecasts and encode the selected regime;
the three smaller checkpoints (4--8B) do neither. The same size-linked split
appears within both Qwen3 and Llama~3.1. Held-out AUC at $k=0$ is $.778$ for
Qwen2.5-32B, $.728$ for Llama~3.1-70B, $.678$ for Qwen3-14B and $.662$ for
Qwen2.5-72B, all rejecting their held-out permutation null at $p<.01$, against
$.566$, $.520$ and $.423$ for Qwen3-8B, Qwen3-4B and Llama~3.1-8B, none of which
reject. Llama~3.1 gives the cleanest contrast: it holds tokenizer and architecture
fixed while varying only scale, and it is the one within-family split whose
behavioral classification also survives the harder population.
Decodability and behavior remain correlational.
Figure~\ref{fig:coin-city-causal-patch} shows an activation-patching test in
Qwen3-14B. Replacing the label-token representation with the representation
induced by the opposite episode-local mapping shifts the zero-shot forecast
toward the donor regime (\(+.333\)). After four target cases, the estimate is
\(-.035\) and lies inside the permutation-null range. At that checkpoint, the
episode-local convention causally contributes to the forecast when behavior
depends on the cue.

\begin{figure}[!ht]
\vspace{-0.4em}
\centering
\includegraphics[width=\linewidth]{exp2_causal_patch.pdf}
\vspace{-0.7em}
\caption{\textbf{Patching the induced label representation changes the zero-shot
forecast, but not the forecast after direct target evidence arrives.} One sealed
\(k{=}0\) episode, selected as the episode whose shift is closest to the
88-episode mean. The donor and recipient prompts are identical except for the two
City C label tokens. Writing the donor's hidden states at those tokens over the
recipient's at layers 19--21, then finishing the forward pass, moves the
recipient's forecast from 51.3 to 54.0 --- the donor's own forecast --- and its
stated rationale switches to the donor's reference city. Writing the recipient's
own states back in leaves the forecast unchanged, and the reverse direction moves
the donor from 54.0 to 51.4. Responsiveness is the implied poll change per news
point. Across-episode means and permutation nulls are in the text and
Table~\ref{tab:coin-city-causal-patch}.}
\label{fig:coin-city-causal-patch}
\vspace{-1.0em}
\end{figure}

\newcommand{\transfergridlines}{11}
% Compact inline orientation graphic for the Experiment 3 transfer design: a 2x2
% grid crossing generative mechanism (rows) with surface domain (columns).
%
% Row cards carry the mechanism detail: a two-node vs three-node graph with the
% mediator drawn hollow (it is never named to the model) and a persistence loop,
% plus the six-period impulse response on labelled axes. Bar heights are the
% normalised responses to a unit shock under worlds.transition at the midpoint of
% each sampling band (direct: gain .57; mediated: gain .57, rho .61, phi .30,
% lambda .81), i.e. A = (1, 0, 0, 0, 0, 0) and B = (.81, .74, .52, .34, .21, .13).
%
% Geometry note: cards have fixed dimensions and every label is a single line
% with no text width, so nothing can rewrap into a neighbour.
% Callers may set \transfergridlines to match the paragraph they anchor this
% to; wrapfig cannot wrap a section heading, so the count must not overrun it.
\providecommand{\transfergridlines}{13}
\begin{wrapfigure}[\transfergridlines]{r}{0.48\textwidth}
\vspace{-0.8em}
\centering
\definecolor{miniCityFill}{HTML}{E8F2FB}
\definecolor{miniCityEdge}{HTML}{4F7EA8}
\definecolor{miniHarborEdge}{HTML}{3F8E82}
\definecolor{miniTestFill}{HTML}{E7F4EA}
\definecolor{miniTestEdge}{HTML}{6FA97F}
\definecolor{miniTestChip}{HTML}{5E9470}
% Coin discs reuse Figure 1's yellow pair exactly -- the fill and stroke of its
% NAIVE box -- rather than the saturated gold used before. Figure 1's other warm
% pair (#FFE6CC on #D79B00, its FREQUENTIST box) is the orange one and is not
% used here. The glyph inside each disc still carries its domain colour.
\definecolor{miniCoinFill}{HTML}{FFF2CC}
\definecolor{miniCoinRim}{HTML}{D6B656}
\definecolor{miniMute}{HTML}{8A93A3}
\definecolor{miniInk}{HTML}{3F4550}
\definecolor{miniBar}{HTML}{7C93AD}
% Coin City: a skyline stamped into a coin.
\newcommand{\coincity}[2]{%
  \begin{scope}[shift={(#1,#2)},scale=0.70]
    \fill[miniCoinFill] (0,0) circle (.37);
    \draw[miniCoinRim,line width=.75pt] (0,0) circle (.37);
    \fill[miniCityEdge] (-.23,-.18) rectangle (-.10,.04);
    \fill[miniCityEdge] (-.07,-.18) rectangle (.07,.16);
    \fill[miniCityEdge] (.10,-.18) rectangle (.24,.09);
    \draw[miniCityEdge,line width=.65pt] (-.27,-.19)--(.27,-.19);
  \end{scope}}
% Coin Harbor: a small vessel and waves stamped into a coin.
\newcommand{\coinharbor}[2]{%
  \begin{scope}[shift={(#1,#2)},scale=0.70]
    \fill[miniCoinFill] (0,0) circle (.37);
    \draw[miniCoinRim,line width=.75pt] (0,0) circle (.37);
    \fill[miniHarborEdge] (-.24,-.05)--(.24,-.05)--(.15,-.18)--(-.15,-.18)--cycle;
    \draw[miniHarborEdge,line width=.7pt]
      (-.01,-.05)--(-.01,.15)--(.15,.02)--(-.01,.02);
    \draw[miniHarborEdge,line width=.55pt]
      (-.26,-.24)..controls(-.18,-.19)and(-.10,-.29)..(-.02,-.24)
      ..controls(.06,-.19)and(.14,-.29)..(.26,-.24);
  \end{scope}}
% Impulse response on labelled axes: response (Delta y) against period (t).
\newcommand{\impulse}[3]{% #1 = axis corner x, #2 = baseline y, #3 = bar heights
  \draw[-{Stealth[length=2.6pt,width=2.2pt]},draw=miniInk!65,line width=.5pt]
    (#1,#2) -- (#1,#2+0.40);
  \draw[-{Stealth[length=2.6pt,width=2.2pt]},draw=miniInk!65,line width=.5pt]
    (#1,#2) -- (#1+0.72,#2);
  \node[anchor=south,font=\sffamily\scriptsize,text=miniInk,inner sep=1pt]
    at (#1,#2+0.40) {$\Delta y$};
  \node[anchor=west,font=\sffamily\scriptsize,text=miniInk,inner sep=1.5pt]
    at (#1+0.72,#2) {$t$};
  \foreach \i/\h in {#3}{%
    \fill[miniBar] (#1+0.055+\i*0.105,#2) rectangle (#1+0.055+\i*0.105+0.070,#2+\h);}}
\resizebox{\linewidth}{!}{%
\begin{tikzpicture}[
  x=1cm, y=1cm,
  cell/.style={rounded corners=3pt, minimum width=1.95cm, minimum height=1.25cm,
    line width=.6pt, inner sep=0pt},
  traincell/.style={cell, draw=miniCityEdge, fill=miniCityFill},
  heldcell/.style={cell, draw=miniTestEdge, fill=miniTestFill,
    dash pattern=on 1.9pt off 1.5pt},
  rowcard/.style={rounded corners=3pt, minimum width=2.60cm,
    minimum height=1.25cm, inner sep=0pt, draw=miniMute!30, fill=black!3,
    line width=.5pt},
  dname/.style={anchor=south, font=\sffamily\bfseries\scriptsize},
  rname/.style={anchor=center, font=\sffamily\scriptsize, text=miniInk},
  note/.style={anchor=center, font=\sffamily\scriptsize, text=miniInk},
  gnode/.style={circle, draw=miniInk!75, fill=white, line width=.5pt,
    inner sep=0pt, minimum size=0.26cm, font=\sffamily\tiny, text=miniInk},
  hidden/.style={gnode, draw=miniMute, fill=miniMute!14,
    dash pattern=on 1.1pt off .9pt},
  garrow/.style={-{Stealth[length=2.4pt,width=2pt]}, draw=miniInk!70,
    line width=.45pt, shorten >=.4pt, shorten <=.4pt},
  chip/.style={anchor=center, rounded corners=1.6pt, draw=none, text=white,
    inner xsep=3pt, inner ysep=1.1pt, font=\sffamily\bfseries\tiny},
]
% ---- domain marks and names ----------------------------------------------
\coincity{0.975}{2.56}
\coinharbor{3.075}{2.56}
\node[dname, text=miniCityEdge]   at (0.975,1.80) {COIN CITY};
\node[dname, text=miniHarborEdge] at (3.075,1.80) {COIN HARBOR};

% ---- Structure A: direct, one period -------------------------------------
\node[rowcard] at (-1.50,1.00) {};
\node[rname] at (-1.50,1.45) {\textbf{Structure A} direct};
\node[gnode] (ax) at (-2.51,0.82) {$x$};
\node[gnode] (ay) at (-2.07,0.82) {$y$};
\draw[garrow] (ax) -- (ay);
\impulse{-1.22}{0.62}{0/0.360, 1/0.000, 2/0.000, 3/0.000, 4/0.000, 5/0.000}

% ---- Structure B: mediated, persistent -----------------------------------
\node[rowcard] at (-1.50,-0.42) {};
\node[rname] at (-1.50,0.03) {\textbf{Structure B} mediated};
\node[gnode]  (bx) at (-2.51,-0.60) {$x$};
\node[hidden] (bm) at (-2.07,-0.60) {$m$};
\node[gnode]  (by) at (-1.63,-0.60) {$y$};
\draw[garrow] (bx) -- (bm);
\draw[garrow] (bm) -- (by);
\draw[garrow] (bm) to[out=118,in=62,looseness=4.5] (bm);
\impulse{-1.22}{-0.80}{0/0.292, 1/0.265, 2/0.188, 3/0.123, 4/0.077, 5/0.048}

% ---- the four evaluation cells -------------------------------------------
\node[traincell] at (0.975,1.00) {};
\node[heldcell]  at (3.075,1.00) {};
\node[heldcell]  at (0.975,-0.42) {};
\node[heldcell]  at (3.075,-0.42) {};

\node[chip, fill=miniCityEdge] at (0.975,1.26) {TRAINED};
\node[chip, fill=miniTestChip] at (3.075,1.26) {TEST};
\node[chip, fill=miniTestChip] at (0.975,-0.16) {TEST};
\node[chip, fill=miniTestChip] at (3.075,-0.16) {TEST};

\node[note] at (0.975,0.80) {in-distribution};
\node[note] at (3.075,0.80) {domain shift};
\node[note] at (0.975,-0.62) {structure shift};
\node[note] at (3.075,-0.62) {joint shift};
\end{tikzpicture}%
}
\vspace{-0.3em}
\caption{\footnotesize\textbf{Transfer grid.} Training uses the shaded cell;
all four are evaluated. Columns vary domain; rows vary response mechanism.}
\label{fig:exp3-transfer-grid}
\vspace{-0.9em}
\end{wrapfigure}


Real markets show that models distinguish probative evidence from topical
similarity. Coin City shows context selecting between relationships visible in
the same prompt, then direct observations revising that choice. Arbitrary labels
rule out a stable semantic prior, and the patch connects one induced
representation to behavior. Together, these results establish relationship
selection and revision within the tested tasks. The next question is whether training carries that
behavior into a new world.

\vspace{3.0\baselineskip}
